\section{Đặt vấn đề}
Trong kỷ nguyên số hoá, lượng tài liệu và tri thức số đã và đang gia tăng một cách chóng mặt. Việc tự học, tự nghiên cứu thông qua tài liệu trực tuyến trở thành một yêu cầu bắt buộc đối với sinh viên và nhà nghiên cứu. Đồng thời, sự phát triển mạnh mẽ của các mô hình ngôn ngữ lớn (LLM) cùng kỹ thuật truy xuất tăng cường tạo sinh (RAG) đã mở ra các hướng đi mới trong việc hỗ trợ người dùng trích xuất thông tin và tương tác với tài liệu. Tuy nhiên, qua quá trình khảo sát hành vi người dùng cùng các giải pháp tồn tại trên thị trường, nhóm nhận thấy còn tồn tại một số điểm nghẽn quan trọng:
\subsection{Sự phân mảnh về trải nghiệm}
Quá trình tự học, tự nghiên cứu đòi hỏi sự tập trung cao độ trong tư duy. Hiện tại, giải pháp được đa số người dùng lựa chọn là sử dụng một ứng dụng để đọc tài liệu (PDF Reader, Kindle) và một ứng dụng khác để tra cứu, hỏi đáp (ChatGPT, Gemini, Perplexity, \ldots). Việc này tạo ra một trải nghiệm gián đoạn, có tính phân mảnh cao: người dùng phải liên tục sao chép nội dung qua lại, tái thiết lập ngữ cảnh cho công cụ AI, đồng thời làm giảm sự tập trung. Sự tách rời giữa luồng đọc tài liệu và luồng tra cứu đặt ra yêu cầu về một không gian làm việc hợp nhất, tích hợp đầy đủ tính năng tương tác với tài liệu và hỏi đáp AI diễn ra đồng thời trên cùng một màn hình.
\subsection{Thao tác thủ công khi cần giải thích một đoạn văn bản}
Khi gặp một đoạn văn bản khó hiểu, người dùng của các công cụ đọc tài liệu phổ thông (không tích hợp AI) phải tự tay sao chép đoạn văn bản, chuyển sang một ứng dụng AI khác, dán vào và gõ thêm câu hỏi. Thao tác rườm rà này làm gián đoạn mạch đọc, và ngữ cảnh gửi cho AI (toàn văn bản hoặc không có gì) thường không khớp chính xác với đoạn văn người dùng thực sự quan tâm, dẫn đến câu trả lời lạc đề hoặc quá chung chung.
\subsection{Thiếu tính minh bạch và khả năng kiểm chứng}
Các công cụ hỏi đáp AI phổ thông thường trả lời mà không chỉ rõ câu trả lời được lấy từ đâu trong tài liệu gốc, khiến người dùng khó kiểm chứng độ chính xác và dễ gặp hiện tượng ảo giác (hallucination) mà không nhận ra. Đối với tài liệu học thuật và kỹ thuật, nơi độ chính xác quan trọng hơn tốc độ, sự thiếu minh bạch này là một rào cản lớn.
\subsection{Thiếu tính cá nhân hoá theo quá trình học tập}
Phần lớn công cụ đọc tài liệu hiện nay xử lý mỗi phiên làm việc độc lập, không ghi nhận và tích lũy các khái niệm mà người dùng đã tìm hiểu qua nhiều tài liệu và nhiều phiên đọc khác nhau. Người dùng không có cách nào để nhìn lại bức tranh tổng thể về những gì bản thân đã học, hay các mối liên hệ giữa các khái niệm rải rác trong nhiều tài liệu.

Từ những điểm nghẽn nêu trên, đồ án hướng tới việc xây dựng DigitalBookLLM, một ứng dụng đọc sách tích hợp AI trong cùng một không gian làm việc. Người dùng có thể bôi đen một đoạn văn bản và nhận câu trả lời bám sát ngữ cảnh đó ngay lập tức, kèm trích dẫn rõ ràng về nguồn gốc thông tin. Đồng thời, hệ thống tự động xây dựng một đồ thị tri thức cá nhân ghi nhận các khái niệm đã tiếp xúc qua quá trình đọc và trò chuyện với AI.
